% Evaluation (Section 5). All numbers come from results/REPORT.md and the logged runs in results/ of the
% SpecLearn artifact.
\section{Evaluation}
\label{sec:eval}\label{sec:evaluation}

This section reports a controlled study of \system{} on four sandboxed domains with reference oracles, which predates the dataset-based formulation of Section~\ref{sec:problem}. % TODO: add the real-world dataset evaluation (docs/EXPERIMENT_PLAN.md)
It answers four research questions:

\begin{description}
  \item[RQ1 (Exact membership).] When membership answers are exact, which equivalence checks let $L^*$ identify $L_G$,
  and how many expert labels does active learning need compared with passive and random-query learning?
  \item[RQ2 (Membership assumption).] How close do LLM membership judges come to the exact answers that assumption
  (A3) requires, and how do their errors propagate into the learned contracts?
  \item[RQ3 (Runtime fidelity).] Do the event abstraction and the compiled AgentSpec contracts enforce the
  learned model faithfully at runtime?
  \item[RQ4 (Safety and utility).] Do the synthesized contracts prevent policy violations of an LLM agent while
  preserving its legitimate tasks, compared with expert-written rules and direct LLM synthesis?
\end{description}

% -------------------------------------------------------------------------------------
\subsection{Experimental Setup}
\label{sec:setup}

\paragraph{Agent domains.}
We built four instrumented agent domains (Table~\ref{tab:domains}). Each provides the inputs of \system{}: a fixed
natural-language policy $P$, a tool API, a deterministic sandbox, and an expert-written event abstraction
$(\Sigma,\Phi,\mu)$ with Boolean recognition predicates over visible fields. Each domain also has a reference DFA and an independent concrete policy checker. The checker inspects the ground-truth
effects of tool calls and never uses $\Sigma$. Together they serve as the ground truth (A5) and as the reference for the
metrics; in the runs below, ground-truth labels are taken from the reference DFA, which agrees with the checker within
the bounds stated next. \dom{DocShare} is the running example of Section~\ref{sec:motivating-example}.

We checked policy sufficiency (Section~\ref{sec:problem}) within bounds: the checker agrees with the reference DFA on the canonical concretization of
every abstract word up to length 4--5 and on every concrete action sequence of length at most 3.

The domains contain the difficulties the method must handle:
\begin{itemize}
  \item unlabeled files (\dom{DocShare});
  \item a multi-effect tool and an outcome-dependent approval step (\dom{Payments});
  \item unannotated SQL migrations (\dom{DevOps});
  \item synonym field names (\dom{DataExport});
  \item alias tools that produce the same effect (three domains).
\end{itemize}
Each domain has 12--13 held-out tasks: 6--7 benign and 6 unsafe, 49 in total. The unsafe tasks ask for an action that
would violate $P$. Each domain also has 8 training tasks, used only by passive inference.

\begin{table}[t]
  \caption{Agent domains. $|\Sigma|$: events; $|Q|$: states of the minimal reference DFA; tasks: held-out
  benign/unsafe tasks.}
  \label{tab:domains}
  \small
  \setlength{\tabcolsep}{3pt}
  \begin{tabular}{@{}p{1.55cm}p{4.0cm}ccc@{}}
    \toprule
    Domain & Policy $P$ (abridged) & $|\Sigma|$ & $|Q|$ & Tasks \\
    \midrule
    \dom{DocShare} & Confidential documents must not become accessible to unauthorized parties. & 3 & 3 & 6/6 \\
    \dom{Payments} & Payments above \$1{,}000 need a manager's approval; an approval covers one execution and is voided by any change. & 4 & 4 & 6/6 \\
    \dom{DevOps} & Deploy only if tests passed after the latest change; a destructive migration needs a fresh backup. & 6 & 5 & 7/6 \\
    \dom{DataExport} & No personal data to external recipients unless anonymized; at most two external reports. & 5 & 6 & 6/6 \\
    \bottomrule
  \end{tabular}
\end{table}

\paragraph{Target agent and LLMs.}
The target agent $A$ is a JSON tool-calling agent built on Qwen3-30B-A3B. It is \emph{not} told the policy, so the
contract has to constrain it independently of its own decisions.

All LLMs are open-weight models served locally with MLX on an Apple M3 Ultra (96\,GB). They use greedy decoding
(temperature 0) and fixed prompts. The roles are:
\begin{itemize}
  \item \emph{Membership judge:} Qwen3-30B-A3B. It is shown the word's canonical execution in the sandbox, a scripted
  concretization rather than an agent run, and uses a step-wise prompt. Its reasoning mode is used only to escalate unresolved queries.
  \item \emph{Scenario writer for agent-realized tests:} Qwen3-14B. Using a separate model lets judge errors
  and tester errors be told apart.
  \item \emph{LLM event recognizer:} Qwen3-30B-A3B.
\end{itemize}
\systemR{} denotes the same configuration with the judge's reasoning mode always enabled.

\paragraph{Implementation.}
\system{} is 5.5k lines of Python.
\begin{itemize}
  \item \emph{Learner.} $L^*$ with Angluin's counterexample processing.
  \item \emph{Membership service.} It caches labels with their provenance and infers the labels implied by prefix
  closure~\cite{hungar2003domain}. It escalates unresolved queries, and it suspends learning if no oracle resolves a query rather than invent a
  label.
  \item \emph{Equivalence check.} The implementation predates the ground-truth formulation of Section~\ref{sec:lstar}:
  it labels equivalence tests with the active membership oracle. With exact oracles these labels are ground truth;
  with an LLM judge they are not. It builds its tests online rather than replaying a fixed dataset. Each equivalence
  query first runs 6 agent-realized tests: the scenario writer turns a target word into a task, the agent executes it
  in a fresh sandbox, and the realized word $\alpha(z)$ is labeled. Up to 120 hypothetical tests follow, from
  hypothesis-guided selection: the access words of the hypothesis states, extended by $\Sigma^{\le 2}$ and the
  distinguishing suffixes $E$, prioritized by transitions that enter or lie next to the rejecting region, then by
  untested event pairs. They predate the use of a fixed labeled dataset.
  \item \emph{Budget.} Runs labeled by an LLM are capped at 400 labels and 12 equivalence queries.
  \item \emph{Compiler.} It emits one AgentSpec rule per rejecting DFA transition, with \texttt{event\_$e$} and
  \texttt{state\_$q$} predicates, and enforces before execution (\texttt{llm\_self\_examine}: the call is not executed, and the agent may continue).
  \item \emph{Runtime recognition.} The runtime uses a hybrid recognizer: Boolean predicates first, and the LLM only
  when they abstain.
\end{itemize}

\paragraph{Compared conditions.}
We compare six conditions:
\begin{enumerate}
  \item \emph{Expert-written AgentSpec rules}, in two variants over a generic past-time predicate library with the
  operators \texttt{ever}, \texttt{last}, \texttt{since} and \texttt{count}. The \emph{careful} rules are correct by construction.
  The \emph{naive} rules contain plausible omissions, such as the naive rules discussed in Section~\ref{sec:motivating-example}.
  \item \emph{LLM direct synthesis.} From $P$, $\Sigma$ and the tool schemas, the LLM writes either a DFA over $\Sigma$
  or AgentSpec rules over the same predicate library, in one call. We report Qwen3-30B-A3B with reasoning, the
  strongest of seven model configurations.
  \item \emph{Passive inference}: RPNI~\cite{oncina1992rpni} on the labeled prefixes of the agent's executions of the training tasks.
  \item \emph{Random-query learning}: RPNI on randomly chosen words, labeled at the label budget of the matched active
  learner.
  \item \emph{$L^*$ with exact oracles}, backed by the reference DFA.
  \item \emph{\system{}} and \emph{\systemR{}} with practical experts. The LLM judge labels both the membership
  queries and the equivalence tests. Section~\ref{sec:problem} argues against this for equivalence; we report it as
  measured; the split teacher of Section~\ref{sec:lstar}, with labels first, remains to be evaluated.
\end{enumerate}

\paragraph{Metrics.}
\begin{itemize}
  \item \emph{Language level} (Definition~\ref{def:policy-language}). The unsafe-permission rate (UP) is
  $|(L(C)\setminus L_G)\cap\Sigma^{\le\ell}|/|\Sigma^{\le\ell}\setminus L_G|$, and the over-restriction rate (OR) is
  $|(L_G\setminus L(C))\cap\Sigma^{\le\ell}|/|L_G\cap\Sigma^{\le\ell}|$. Unless stated otherwise, $\ell=6$ ($\ell=5$ for
  \dom{DevOps}). Each word is replayed through the contract's monitor, so DFA-based and rule-based contracts are measured
  identically. The rates depend on $\ell$: \systemR{}'s over-restriction on \dom{DevOps} is 5.2\% at $\ell=5$ but
  19.0\% at $\ell=10$.
  \item \emph{Online.} The agent runs every held-out task with the contract in the loop. We count the unsafe tasks in
  which a violation is committed and the benign tasks completed without violation.
  \item \emph{Cost.} Expert labels and LLM tokens.
\end{itemize}

\paragraph{Reproducibility.}
Randomness is seeded. Every LLM response is cached under a hash of the model revision, prompt and decoding parameters.
Re-running all experiments from the cache reproduced all 191 result files exactly.
% TODO: anonymized artifact link

% -------------------------------------------------------------------------------------
\subsection{RQ1: Learning with Exact Membership}
\label{sec:rq1}

\paragraph{Controlled targets.}
We first validated the learner on 200 random minimal prefix-closed DFAs with 4 to 64 states and 2 to 16 events,
10 per combination, each learned under 15 configurations (3{,}000 runs). We used an exact membership oracle, an exact equivalence oracle that returns the
shortest counterexample or one padded to a length of at least 10, 25, 50 or 100, and three counterexample-processing strategies: Angluin's~\cite{angluin1987learning}, Maler--Pnueli's~\cite{maler1995learnability},
and Rivest--Schapire's~\cite{rivest1993inference}.

\begin{itemize}
  \item $L^*$ recovered the minimal target in all 3{,}000 runs and never needed more than $N-1$ counterexamples.
  \item The number of distinct membership queries grew roughly as $N^{1.6}$ and as $|\Sigma|^{1.1}$. It stayed below
  49\% (median 8\%) of the classical bound $O(|\Sigma|(B+1)N^2)$.
  \item Long counterexamples inflate the queries of Angluin's strategy ($\times 4.9$ at $B=100$, $N=32$, $|\Sigma|=4$)
  and of Maler--Pnueli's ($\times 7.9$), but not those of Rivest--Schapire's (1{,}606 vs.\ 1{,}685).
  \item On \dom{DocShare} the observation table reproduces the walk-through of Section~\ref{sec:motivating-example}. The first hypothesis has one
  state; the verified counterexample \texttt{c\,p} adds suffix \texttt{p} through a consistency fix; MQ(\texttt{c\,n\,p})
  $=1$ separates \texttt{c} from \texttt{c\,n}; the result is the three-state DFA of Figure~\ref{fig:sharing-dfa}.
\end{itemize}

\paragraph{Domains.}
Table~\ref{tab:rq1} shows the label cost. With exact labels, the equivalence tests are ground-truth-labeled: agent-realized tests
form an online corpus of labeled agent trajectories, and the guided suite adds labeled abstract words. Neither is
checked for completeness. % TODO: replace with the real-world dataset experiments of docs/EXPERIMENT_PLAN.md
\begin{itemize}
  \item With exact equivalence queries, $L^*$ recovers all four reference DFAs with 11--111 labels.
  \item Agent-realized tests alone, with exact labels and 10 tests per equivalence query (instead of the default 6), recover \dom{DocShare} and
  \dom{DevOps} but stop prematurely on \dom{Payments} and \dom{DataExport}, leaving 4.3\% and 17.4\% unsafe
  permissions. The missing distinctions need long sequences (approval consumption, \texttt{ml\,ap\,ex\,ex}; a third
  external send). Every disagreement this corpus exposed was genuine, but it ran out of disagreements before the hypothesis was correct:
  the corpus was incomplete.
  \item The hypothesis-guided suite recovers all four, at 27--311 labels. Its final query spends its budget without finding
  a disagreement, so the result is only consistent with its tests; nothing certifies it.
  \item Random-word equivalence testing recovers \dom{DataExport} in only 4 of 10 seeds.
  \item At the same label budgets, random-query learning recovers only \dom{DocShare} (0 of 10 seeds elsewhere, up to
  7.1\% UP).
  \item Passive inference from the training traces recovers only \dom{DocShare} (up to 42\% UP elsewhere).
\end{itemize}

We also swept the equivalence budget (exact labels, 20 seeds). Hypothesis-guided testing removes premature stops at 40 labels
per equivalence query on \dom{Payments} and \dom{DevOps}, and at 160 on \dom{DataExport}. At 160 labels, random
testing still stops prematurely in 5\% (\dom{Payments}) and 50\% (\dom{DataExport}) of runs. Random testing is
competitive only where violations are dense: on \dom{DevOps} with 10 labels it stops prematurely in 10\% of runs, against
80\% for hypothesis-guided testing.

\begin{table}[t]
  \caption{Learning with exact labels. $L^*$ and Guided EQ (hypothesis-guided bounded testing): expert labels until
  termination (both learned the exact DFA). Random EQ and Random query: seeds (of 10) that learned $L_G$ exactly; Random
  query uses the Guided-EQ label budget. Passive: labels used and whether the result is exact.}
  \label{tab:rq1}
  \small
  \setlength{\tabcolsep}{3pt}
  \begin{tabular}{@{}lccccc@{}}
    \toprule
    Domain & $L^*$ & Guided & Random & Random & Passive \\
    & (exact EQ) & EQ & EQ & query & \\
    \midrule
    \dom{DocShare}   & 11  & 27  & 10/10 & 10/10 & 10, yes \\
    \dom{Payments}     & 36  & 103 & 10/10 & 0/10  & 14, no \\
    \dom{DevOps}       & 73  & 195 & 10/10 & 0/10  & 14, no \\
    \dom{DataExport} & 111 & 311 & 4/10  & 0/10  & 10, no \\
    \bottomrule
  \end{tabular}
\end{table}

\answerbox{RQ1}{With exact membership and equivalence oracles, $L^*$ identifies every target, as classical $L^*$
guarantees: all 3{,}000 synthetic runs (200 targets) and all four domain policies. With ground-truth labels
but no completeness, agent-realized trajectories alone stop early on two of four domains, and the budget-capped
hypothesis-guided suite recovers all four without a certificate. Learning against a fixed real-world labeled dataset is not yet evaluated.
Random-query learning at the same label budget and passive learning from the training traces each learn only one of
the four contracts.}

% -------------------------------------------------------------------------------------
\subsection{RQ2: The Membership Assumption}
\label{sec:rq2}

\paragraph{Membership judges.}
We measured judges on 178 independently labeled words (Table~\ref{tab:judges}). The set contains all or sampled words of
length at most 2, plus 24 longer words per domain balanced over three categories: permitted, violating at the last event,
and violating earlier.

\begin{itemize}
  \item The primary judge is 87.6\% accurate (100/87/81/85\% per domain).
  \item Abstract evidence gives the same accuracy as grounded evidence but shifts the errors: more over-rejection
  (19.2\% vs.\ 13.5\%) and fewer unsafe acceptances (1.4\% vs.\ 8.1\%).
  \item A direct (non-step-wise) prompt or a paraphrased yes/no prompt costs 7 points.
  \item With reasoning enabled, the judge labels all 178 words correctly, at about $1.8\times$ the tokens.
  \item Verdicts are deterministic: 48 judgments regenerated without the cache were all reproduced. They are not robust
  to wording: the primary judge changes its verdict on 15--19\% of words under a paraphrased prompt, renamed event
  symbols or abstract evidence.
  \item Other models fail in opposite directions. Llama-3.1-8B accepts 69\% of prohibited words; Qwen3-4B and Qwen3-14B
  reject about half of the permitted ones.
\end{itemize}

\begin{table}[t]
  \caption{LLM membership judges on 178 labeled words. Acc.: accuracy; UA: prohibited words judged permitted; OR:
  permitted words judged violating.}
  \label{tab:judges}
  \small
  \setlength{\tabcolsep}{4pt}
  \begin{tabular}{@{}lrrr@{}}
    \toprule
    Judge configuration & Acc. & UA & OR \\
    \midrule
    Qwen3-30B, grounded, step-wise (primary) & 87.6 & 8.1 & 13.5 \\
    \quad abstract evidence                  & 87.6 & 1.4 & 19.2 \\
    \quad direct prompt                      & 80.9 & 10.8 & 19.2 \\
    \quad paraphrased (yes/no) prompt        & 80.3 & 9.5 & 25.0 \\
    \quad reasoning enabled                  & \textbf{100.0} & 0.0 & 0.0 \\
    Qwen3-32B, grounded, direct              & 91.0 & 10.8 & 7.7 \\
    Mistral-Small-24B, grounded, direct      & 86.0 & 1.4 & 23.1 \\
    Llama-3.1-8B, grounded, direct           & 70.8 & 68.9 & 1.0 \\
    Qwen3-4B, grounded, direct               & 68.5 & 10.8 & 46.2 \\
    Qwen3-14B, grounded, direct              & 68.0 & 0.0 & 51.9 \\
    \bottomrule
  \end{tabular}
\end{table}

\paragraph{Sensitivity to label errors.}
With simulated persistent noise, each word's label is flipped with probability $\rho$ (20 seeds, hypothesis-guided
equivalence testing). Exact recovery falls from 100\% to 5--75\% at $\rho=1\%$, and to 0--15\% at $\rho=5\%$. At
$\rho=5\%$ the contracts over-restrict 23--29\% of permitted words and permit 5--9\% of prohibited ones (words up to
length 8).

\paragraph{Learning with LLM judges.}
Table~\ref{tab:rq2} shows the learned contracts. Both configurations learn \dom{DocShare} exactly from 29 labels.

Without reasoning, the judge mislabels 19--36\% of the queries $L^*$ asks in the other domains. The errors are mostly
over-cautious ``violation'' verdicts, for instance treating setting a large draft payment as paying it, plus lost counts
on long sequences. $L^*$ exhausts its budget with 9--17-state monitors that over-restrict 15--61\% of permitted words.

With reasoning, the error rate drops to 0.3--6\%. The \dom{Payments} contract is practically exact (one over-restricted
word out of 3{,}910), and \dom{DevOps} reaches 3.3\% UP and 5.2\% OR. Three effects remain:
\begin{enumerate}
  \item \emph{One wrong label can prevent convergence.} On \dom{Payments}, a single wrong label kept $L^*$ from
  converging: it spent the remaining budget adding states to explain it (10 states instead of 4). Because the same
  judge labeled the equivalence tests, they could not expose the error; a label for that word would have overridden it.
  \item \emph{Balanced word sets overstate accuracy.} The judge that was perfect on the balanced word set errs on
  3.5--6\% of the learner's queries, which are long and probe the policy boundary. On \dom{DataExport} it consistently
  judged reading personal data and then sending it to an internal recipient to be a violation.
  \item \emph{Long queries can exhaust the reasoning budget.} Learning on \dom{DataExport} was suspended when the judge
  used up its reasoning budget on a 10-event query without reaching a verdict.
\end{enumerate}

\begin{table}[t]
  \caption{Contracts learned with LLM judges. States: learned/reference. Wrong: mislabeled queries / queries judged. UP/OR (\%): language level,
  all words up to length 6 (5 for \dom{DevOps}). $^\dagger$label budget exhausted; $^\ddagger$suspended on an unresolved query.}
  \label{tab:rq2}
  \small
  \setlength{\tabcolsep}{3pt}
  \begin{tabular}{@{}lcrr@{\hspace{8pt}}crr@{}}
    \toprule
    & \multicolumn{3}{c}{\system{}} & \multicolumn{3}{c}{\systemR{}} \\
    \cmidrule(r){2-4}\cmidrule(l){5-7}
    Domain & States & Wrong & UP/OR & States & Wrong & UP/OR \\
    \midrule
    \dom{DocShare}   & 3/3 & 0/29 & 0.0/0.0 & 3/3 & 0/29 & 0.0/0.0 \\
    \dom{Payments}     & 13/4$^\dagger$ & 144/400 & 17.2/60.2 & 10/4$^\dagger$ & 1/399 & 0.0/0.03 \\
    \dom{DevOps}       & 17/5$^\dagger$ & 94/395  & 4.7/60.5  & 13/5$^\dagger$ & 14/399 & 3.3/5.2 \\
    \dom{DataExport} & 9/6$^\dagger$  & 74/399  & 7.4/15.2  & 11/6$^\ddagger$ & 16/266 & 17.3/10.2 \\
    \bottomrule
  \end{tabular}
\end{table}

\paragraph{Agent-realized tests.}
The scenario writer and the agent reproduced the target word exactly in 53--94\% of tests. Within \system{},
agent-realized tests produced 1--3 verified counterexamples per domain. Used alone with exact labels, they stop early
on two domains (Section~\ref{sec:rq1}).

\answerbox{RQ2}{Current LLM judges do not meet assumption (A3), and $L^*$ amplifies even a few wrong labels into
oversized, over-restrictive monitors. A non-reasoning judge (88\% accurate) yields one exact contract out of four; the same model
with reasoning yields one exact and one practically exact contract. Judge accuracy must be measured on the learner's own
queries, which are harder than balanced samples.}

\begin{table*}[t]
  \caption{Safety and utility. Language level: unsafe-permission / over-restriction rate (\%) over all words up to length
  6 (5 for \dom{DevOps}). Online: unsafe held-out tasks with a committed violation, and benign tasks completed without
  violation. LLM direct synthesis uses Qwen3-30B-A3B with reasoning. $^\dagger$Labeled by the same LLM judge as
  \system{}: random-query learning at \system{}'s label budget, passive inference on the labeled training traces. Online runs were performed for the listed conditions only.}
  \label{tab:rq4}
  \small
  \begin{tabular}{@{}lccccrr@{}}
    \toprule
    & \multicolumn{4}{c}{Language level: UP / OR (\%)} & \multicolumn{2}{c}{Online} \\
    \cmidrule(r){2-5}\cmidrule(l){6-7}
    Condition & \dom{DocShare} & \dom{Payments} & \dom{DevOps} & \dom{DataExport} & Violations & Benign done \\
    \midrule
    No contract                        & 100 / 0.0 & 100 / 0.0 & 100 / 0.0 & 100 / 0.0 & 21/24 & 24/25 \\
    Expert rules, careful              & 0.0 / 0.0 & 0.0 / 0.0 & 0.0 / 0.0 & 0.0 / 0.0 & 0/24 & 23/25 \\
    Expert rules, naive                & 0.0 / 26.4 & 40.4 / 0.0 & 14.6 / 0.0 & 16.3 / 8.8 & 7/24 & 19/25 \\
    LLM direct synthesis, DFA          & 0.0 / 0.0 & 4.3 / 0.0 & 0.0 / 4.0 & 0.0 / 0.0 & 2/24 & 22/25 \\
    LLM direct synthesis, rules        & 0.0 / 0.0 & 4.3 / 0.0 & 28.5 / 28.7 & 0.0 / 0.0 & 4/24 & 23/25 \\
    Passive inference (RPNI)$^\dagger$ & 0.0 / 0.0 & 10.6 / 34.3 & 14.0 / 35.3 & 42.4 / 27.9 & 3/24 & 20/25 \\
    Random-query learning$^\dagger$    & 0.0 / 0.0 & 5.7 / 43.7 & 6.7 / 44.7 & 25.7 / 11.1 & 1/24 & 21/25 \\
    $L^*$ with exact oracles           & 0.0 / 0.0 & 0.0 / 0.0 & 0.0 / 0.0 & 0.0 / 0.0 & 0/24 & 22/25 \\
    \midrule
    \system{}                       & 0.0 / 0.0 & 17.2 / 60.2 & 4.7 / 60.5 & 7.4 / 15.2 & 3/24 & 19/25 \\
    \systemR{}                      & 0.0 / 0.0 & 0.0 / 0.03 & 3.3 / 5.2 & 17.3 / 10.2 & 1/24 & 22/25 \\
    \bottomrule
  \end{tabular}
\end{table*}

% -------------------------------------------------------------------------------------
\subsection{RQ3: Runtime Fidelity}
\label{sec:rq3}

\paragraph{Event recognition.}
Table~\ref{tab:recognizers} summarizes recognition on the scripted corpus.
\begin{itemize}
  \item \emph{Boolean predicates} never misclassify an effect, but they abstain on up to 40\% of effects, such as
  unlabeled files, unannotated SQL and synonym field names. The conservative monitor turns abstentions into
  over-blocking.
  \item \emph{LLM predicates} never abstain, but misclassify 20--32\% of effects. With the \emph{reference} DFA, this
  alone causes 31 false permissions.
  \item \emph{Hybrid recognition} is 93--100\% accurate with zero false permissions. Its residual over-blocks come from
  the LLM misreading an unannotated \texttt{UPDATE} migration or an unlabeled public file.
\end{itemize}

\begin{table}[t]
  \caption{Event recognition on 300 scripted executions per domain: accuracy (abstentions) in \%, and the false
  permissions (FP) and over-restrictions (OR) that recognition alone causes under the reference DFA, summed over domains.}
  \label{tab:recognizers}
  \small
  \setlength{\tabcolsep}{3.5pt}
  \begin{tabular}{@{}lccccrr@{}}
    \toprule
    Recognizer & \dom{Doc} & \dom{Pay} & \dom{Dev} & \dom{Data} & FP & OR \\
    \midrule
    Boolean & 60 (40) & 100 (0) & 87 (13) & 98 (2) & 0 & 102 \\
    LLM     & 80 (0)  & 68 (0)  & 73 (0)  & 78 (0) & 31 & 136 \\
    Hybrid  & 98 (0)  & 100 (0) & 93 (0)  & 100 (0) & 0 & 51 \\
    \bottomrule
  \end{tabular}
\end{table}

\paragraph{Policy sufficiency.}
To find collisions, we realized every abstract word up to length 3--4 through all concrete call sequences of the sandbox
and labeled each realization with the independent checker.
\begin{itemize}
  \item With exact Boolean recognition, the expert alphabet produces no contradictory word in any domain (873 realizable
  words).
  \item With the hybrid recognizer used at runtime, LLM misreadings make 24.2\% (\dom{DocShare}) and 16.7\%
  (\dom{DevOps}) of the realizable words contradictory: recognition errors break sufficiency (A1) just as a coarse
  alphabet does.
  \item A naive ``event = tool name'' alphabet has 19--63\% contradictory words.
  \item Vocabularies derived by two LLMs (ablation~i) drop the distinctions that matter: file sensitivity, the payment
  threshold and approval semantics, and destructiveness. They are contradictory on up to 69\% of words (\dom{DevOps}).
  Contracts learned over them with an exact grounded oracle give 7--100\% false permissions.
\end{itemize}

\paragraph{Compilation.}
All 142 DFA-based contracts produced in the evaluation agree with the DFA they were compiled from on every reachable
transition (4{,}059 decision and state-update checks) and on all 1.37\,M words up to the length bound. Executed end to end
on concrete calls, including alias tools, the reference and exact-$L^*$ contracts agree with their DFA on 100\% of
decisions.

One contract deviates: the noisy \system{} contract for \dom{DevOps}, at 89\%. The cause is semantic, not the
compiler. Its DFA rejects one outcome of \texttt{run\_tests}, and the outcome is unknown before execution, so the
pre-effect check must block the call.

\paragraph{Instrumentation and timing.}
Leaving one alias tool out of $\mu$ lowers trigger interception recall from 100\% to 18--94\%. Online, the agent then
bypassed the contract through the unmapped tool in 2--3 of the 6 unsafe episodes per domain (33--50\% violations
instead of 0\%). Used as a post-effect detector, the same contract flags violations but prevents none: 67--100\% of
unsafe episodes violate. A monitor check takes 1.5--2.1\,$\mu$s with Boolean recognition.

\answerbox{RQ3}{Compilation is faithful for all 142 contracts. The weak points lie at the boundary between the model and
the agent:
\begin{itemize}
  \item recognizers must abstain rather than guess, and hybrid recognition eliminates the false permissions;
  \item the alphabet must preserve the policy's distinctions, which the expert alphabet does and LLM-derived ones do
  not, and recognition must not merge them, which the deployed hybrid recognizer does on 17--24\% of words in two
  domains;
  \item every tool that produces a relevant effect must be hooked;
  \item enforcement must happen before the effect commits.
\end{itemize}}

% -------------------------------------------------------------------------------------
\subsection{RQ4: Safety and Utility}
\label{sec:rq4}

Table~\ref{tab:rq4} compares all conditions.

\paragraph{Baselines.}
\begin{itemize}
  \item Without a contract, the agent commits a violation in 21 of the 24 unsafe held-out tasks.
  \item Contracts learned from exact labels and the careful expert rules are exact in all four domains. They prevent
  every violation online while completing 22--23 of the 25 benign tasks. The two benign failures in \dom{DocShare} are
  shared by all contracts: the hybrid recognizer classifies an unlabeled public file as confidential.
  \item The naive expert rules miss approval consumption, the test-after-change requirement and the egress limit
  (15--40\% UP), giving violations in 7 tasks.
  \item Direct synthesis with reasoning is exact on \dom{DocShare} and \dom{DataExport}. It misses approval consumption
  on \dom{Payments} (4.3\% UP), and its rules leave 28.5\% UP on \dom{DevOps}, which leads to 2 (DFA) and 4 (rules)
  online violations. Without reasoning, the same model produces monitors with up to 100\% UP or 94\% OR.
\end{itemize}

\paragraph{\system{}.}
With the non-reasoning judge, \system{} lets 3 violations through and completes 19 benign tasks. \systemR{} lets 1
violation through and completes 22 benign tasks, the same utility as the exact contracts.

\paragraph{Cost.}
Exact $L^*$ needs 11--111 labels (27--311 with bounded equivalence testing). With an LLM judge, \system{} spent 29
labels on \dom{DocShare} and its full 400-label budget elsewhere; \systemR{} was suspended after 266 labels on
\dom{DataExport}. That is 45k--480k LLM tokens per contract without
reasoning and 56k--972k with reasoning. Direct synthesis takes one call of 1.6k--8.1k tokens.


\answerbox{RQ4}{Contracts learned with exact oracles prevent every violation, as the careful expert rules do, and
complete 22 of 25 benign tasks (the careful rules complete 23), without any manually written rule. With practical LLM
judges, \systemR{} lets 1 of 24 unsafe tasks through and completes 22 benign tasks. Online, this is on par with direct
LLM synthesis (2--4 of 24) and random-query learning with the same judge (1 of 24): 24 tasks are too few to separate
them, and at the language level direct DFA synthesis is more accurate than \systemR{} on \dom{DevOps} and
\dom{DataExport}. All are well ahead of naive expert rules (7 of 24). \systemR{} costs up to about 1M tokens per
contract.}

% -------------------------------------------------------------------------------------
\subsection{Threats to Validity}
\label{sec:threats}

\paragraph{Internal.}
We wrote the domains, the reference DFAs, the concrete checkers, the expert abstraction and both expert rule sets. The
careful rules are correct by construction; the naive rules are illustrative and do not estimate how often developers
make such mistakes.

\paragraph{External.}
The four domains are small sandboxes ($|\Sigma|\le 6$, at most 6 states). All LLMs are local open-weight models rather
than frontier models. No human experts took part: simulated label noise stands in for them. AgentSpec contracts run on a
re-implemented runtime with AgentSpec's rule syntax and semantics.

\paragraph{Construct.}
Language-level rates weight all words up to the length bound equally, while online measures weight the behaviors the
agent actually produces; we report both. With 49 held-out tasks, online rates are uncertain. LLM-labeled runs are capped at
400 labels, so runs that do not converge are reported with their last hypothesis.
